\chapter{引言}\label{chap:intro}

\section{研究背景}\label{sec:background}

期权是金融衍生品市场的核心工具，广泛用于风险对冲与投机套利。自Black、Scholes与Merton于1973年提出BSM模型\cite{black1973pricing,merton1973theory}以来，期权定价理论经历了从常数波动率到局部波动率、再到随机波动率的演进。BSM模型具有封闭解析解，适用于短期定价，但无法解释市场中普遍存在的波动率微笑现象。Cox和Ross\cite{cox1976valuation}提出的CEV模型将波动率推广为依赖于标的价格的局部波动率，能够捕捉股票市场的杠杆效应。Heston\cite{heston1993closed}进一步引入随机波动率过程，更准确地描述波动率微笑和厚尾收益率分布，但其二维PDE的数值求解代价显著高于一维情形。

传统数值方法（有限差分、有限元）在求解高维PDE时面临维数灾难：网格点数随维度指数增长，Heston这样的二维模型已需要大量计算资源，而更高维的随机波动率模型几乎无法实时求解。此外，传统工具要求用户手动指定模型类型和参数，对非专业用户门槛较高。

物理信息神经网络（PINN）\cite{raissi2019physics}将PDE残差直接嵌入损失函数，无需网格剖分，训练完成后可对任意输入即时推断（单次推断约2ms）。近年来，PINN已被应用于BSM\cite{dhiman2023pinn}、CEV\cite{zamxaka2023cev}、Heston\cite{hainaut2024heston,guo2024icpinn}等期权定价模型。与此同时，大语言模型（LLM）在自然语言理解和结构化信息提取方面取得了突破性进展\cite{liu2025pdeagent}，为构建自然语言驱动的科学计算接口提供了基础。

\section{现有工作的局限性}\label{sec:motivation}

现有PINN期权定价研究存在两方面主要不足。

\textbf{模型孤立。}现有工作（Dhiman等\cite{dhiman2023pinn}、Hainaut等\cite{hainaut2024heston}）均针对单一模型分别训练独立的PINN，三个模型之间无法共享网络权重，当需要同时支持多种模型时，必须维护多套独立的训练流程和模型文件。

\textbf{交互门槛高。}现有PINN工具要求用户手动指定模型类型和所有数值参数。选择合适的定价模型本身需要量化金融背景，普通用户难以判断在当前市场条件下应使用BSM、CEV还是Heston。

\section{研究目标与贡献}\label{sec:contributions}

本文研究基于物理信息神经网络（Physics-Informed Neural Network，PINN）的多模型期权定价方法，通过统一PDE算子将BSM、CEV、Heston三大模型纳入单一网络，并以大语言模型（Large Language Model，LLM）作为智能路由层实现自然语言交互。在方法设计上，本文提出软掩码机制$\text{mask}=\tanh(\xi/0.05)^2$，将三大模型的扩散系数统一为连续参数化形式，避免了硬切换带来的梯度不连续；采用加法参数化输出$\hat{V}=V_\text{BS}+K\cdot\text{net}(x)$，以Black-Scholes解析解为基线，从根本上解决了深度虚值区域的数值退化问题，同时为网络提供了量级正确的训练起点；引入DeepSeek LLM作为路由层，通过结构化JSON输出实现可靠的参数提取与模型选择，支持中英文自然语言输入。实验表明，三种模型的平值附近相对误差均低于1\%，在BSM和CEV上的精度优于独立训练的PINN基线，LLM路由层在全部测试用例上实现了100\%的模型选择准确率，消融实验进一步验证了各核心组件的必要性。

\section{论文结构}\label{sec:structure}

第二章回顾三大期权定价模型的PDE形式、PINN方法的基本原理及其在期权定价中的应用，以及LLM辅助科学计算的进展。第三章介绍多模型PINN方法，包括统一PDE算子、软掩码机制、加法参数化输出、网络架构、损失函数、训练策略和LLM路由层。第四章报告实验结果，包括精度评估、消融实验和真实SPY期权市场回测。第五章总结全文并讨论局限性与未来工作。
